\documentclass[11pt,a4paper]{article}

% ============================================================
% Arquitectura de despliegue — Portal Agencia ITRC
% ------------------------------------------------------------
% Compilar con:
%   pdflatex arquitectura-portal-itrc.tex
%   pdflatex arquitectura-portal-itrc.tex   (segunda pasada para TOC)
% Requiere: texlive-latex-base, texlive-fonts-recommended,
%           texlive-pictures (para TikZ), texlive-lang-spanish
% ============================================================

\usepackage[utf8]{inputenc}
\usepackage[T1]{fontenc}
\usepackage{lmodern}  % fuentes escalables (requerido por microtype)
\usepackage[spanish,es-tabla]{babel}
\usepackage[margin=2.2cm]{geometry}
\usepackage{microtype}
\usepackage{hyperref}
\usepackage{xcolor}
\usepackage{tikz}
\usepackage{tcolorbox}
\usepackage{booktabs}
\usepackage{fancyhdr}
\usepackage{titlesec}
\usepackage{enumitem}

\usetikzlibrary{
  positioning,
  shapes.geometric,
  shapes.misc,
  arrows.meta,
  fit,
  backgrounds,
  calc
}

% Colores institucionales ITRC
\definecolor{itrcNavy}{HTML}{002147}
\definecolor{itrcGold}{HTML}{B38B40}
\definecolor{govBlue}{HTML}{0943B5}
\definecolor{softGray}{HTML}{F5F5F5}
\definecolor{midGray}{HTML}{888888}
\definecolor{accentGreen}{HTML}{2E7D32}
\definecolor{warnRed}{HTML}{C62828}

% Hyperref
\hypersetup{
  colorlinks=true,
  linkcolor=itrcNavy,
  urlcolor=govBlue,
  citecolor=itrcGold,
  pdftitle={Arquitectura de despliegue — Portal ITRC},
  pdfauthor={Equipo técnico Portal ITRC},
  pdfsubject={Documento técnico de arquitectura},
  pdfkeywords={ITRC, Astro, JAMstack, nginx, GitHub, Sveltia CMS}
}

% Títulos institucionales
\titleformat{\section}
  {\color{itrcNavy}\Large\bfseries}{\thesection}{0.8em}{}
  [\vspace{-0.3em}\color{itrcGold}\rule{\linewidth}{1pt}]
\titleformat{\subsection}
  {\color{itrcNavy}\large\bfseries}{\thesubsection}{0.7em}{}
\titleformat{\subsubsection}
  {\color{itrcNavy}\normalsize\bfseries}{\thesubsubsection}{0.6em}{}

% Header/Footer
\pagestyle{fancy}
\fancyhf{}
\fancyhead[L]{\small\color{itrcNavy}\textbf{Portal ITRC}}
\fancyhead[R]{\small\color{midGray}Arquitectura de despliegue}
\fancyfoot[C]{\small\thepage}
\renewcommand{\headrulewidth}{0.4pt}
\renewcommand{\footrulewidth}{0pt}

% Estilo callouts
\newtcolorbox{callout}[1][]{
  colback=softGray,
  colframe=itrcGold,
  boxrule=1pt,
  arc=2pt,
  left=8pt, right=8pt, top=6pt, bottom=6pt,
  fonttitle=\bfseries,
  #1
}
\newtcolorbox{warnbox}[1][]{
  colback=softGray,
  colframe=warnRed,
  boxrule=1pt,
  arc=2pt,
  left=8pt, right=8pt, top=6pt, bottom=6pt,
  fonttitle=\bfseries,
  #1
}

% Estilos TikZ reutilizables
\tikzset{
  nodoPrincipal/.style={
    rectangle,
    rounded corners=4pt,
    draw=itrcNavy,
    fill=white,
    line width=1pt,
    minimum width=3.4cm,
    minimum height=1.1cm,
    align=center,
    font=\small\bfseries,
    text=itrcNavy
  },
  nodoServidor/.style={
    nodoPrincipal,
    fill=itrcNavy,
    text=white
  },
  nodoGithub/.style={
    nodoPrincipal,
    fill=itrcGold!25,
    draw=itrcGold
  },
  nodoUsuario/.style={
    rectangle,
    rounded corners=12pt,
    draw=govBlue,
    fill=govBlue!15,
    line width=1pt,
    minimum width=2.8cm,
    minimum height=0.9cm,
    align=center,
    font=\small\bfseries,
    text=govBlue
  },
  nodoExterno/.style={
    rectangle,
    rounded corners=2pt,
    draw=midGray,
    fill=softGray,
    line width=0.8pt,
    dashed,
    minimum width=3cm,
    minimum height=0.9cm,
    align=center,
    font=\footnotesize,
    text=midGray
  },
  flecha/.style={
    -{Stealth[length=3mm]},
    line width=0.9pt,
    draw=itrcNavy
  },
  flechaRoja/.style={
    flecha,
    draw=warnRed
  },
  flechaVerde/.style={
    flecha,
    draw=accentGreen
  },
  etiqueta/.style={
    font=\scriptsize,
    text=midGray,
    fill=white,
    inner sep=1pt
  }
}

% ============================================================
\begin{document}

% Portada
\begin{titlepage}
  \centering
  \vspace*{2cm}
  {\color{itrcGold}\rule{\linewidth}{2pt}}\\[0.4cm]
  {\Huge\bfseries\color{itrcNavy} Arquitectura de despliegue\\[0.3cm]}
  {\LARGE\color{itrcNavy} Portal institucional}\\[0.3cm]
  {\Large\color{itrcNavy} Agencia ITRC}\\[0.3cm]
  {\color{itrcGold}\rule{\linewidth}{2pt}}\\[1.5cm]

  {\large Documento técnico de referencia}\\[0.3cm]
  {\small Dirigido a equipo de infraestructura y administradores del sitio}\\[3cm]

  \begin{tcolorbox}[
    colback=softGray,
    colframe=itrcNavy,
    boxrule=1pt,
    width=0.75\linewidth,
    arc=3pt
  ]
  \textbf{Contenido del documento:}
  \begin{itemize}[leftmargin=*, noitemsep, topsep=2pt]
    \item Migración de WordPress a JAMstack (Astro + Sveltia CMS)
    \item Arquitectura de despliegue sobre servidor propio (Ubuntu + nginx)
    \item Flujo de trabajo Git con múltiples webmasters
    \item Autenticación y administración de usuarios
    \item Manejo de contenido estático y binarios
  \end{itemize}
  \end{tcolorbox}

  \vfill
  {\small\color{midGray} Versión 1.0 \textbar{} 2026}
\end{titlepage}

\tableofcontents
\newpage

% ============================================================
\section{Resumen ejecutivo}

El portal institucional de la Agencia ITRC migró desde una arquitectura WordPress (dinámica, basada en base de datos MySQL y PHP) hacia una arquitectura \textbf{JAMstack} basada en el generador de sitios estáticos \href{https://astro.build}{Astro} y el gestor de contenido \href{https://github.com/sveltia/sveltia-cms}{Sveltia CMS}. El portal se despliega sobre el servidor institucional del ITRC (Ubuntu con nginx).

\begin{callout}[title=Cambios clave respecto al modelo anterior]
\begin{itemize}[leftmargin=*, noitemsep]
  \item \textbf{Sin base de datos}: todo el contenido es archivos JSON y Markdown versionados en Git.
  \item \textbf{Sin PHP}: el servidor únicamente sirve archivos estáticos (HTML, CSS, JS, PDF, imágenes).
  \item \textbf{Git como fuente de verdad}: cada cambio es un \textit{commit} con trazabilidad y posibilidad de \textit{rollback}.
  \item \textbf{Despliegue automatizado}: cada \texttt{git push} a la rama principal dispara la publicación en el servidor.
  \item \textbf{Múltiples webmasters}: cada uno con cuenta GitHub propia y acceso auditado.
\end{itemize}
\end{callout}

\subsection{Componentes principales}

\begin{center}
\begin{tabular}{@{}lp{10cm}@{}}
\toprule
\textbf{Componente} & \textbf{Rol} \\
\midrule
\textbf{Astro} & Generador de sitio estático. Transforma archivos JSON/MD en HTML. \\
\textbf{Sveltia CMS} & Panel de administración visual (editor de contenido) que commitea al repo Git. \\
\textbf{GitHub} & Almacenamiento del código, control de versiones, y orquestador de despliegues. \\
\textbf{GitHub Actions} & Pipeline CI/CD que construye y sube el sitio al servidor. \\
\textbf{Ubuntu + nginx} & Servidor institucional que sirve los archivos al usuario final. \\
\textbf{OAuth Proxy} & Intermediario que permite al CMS autenticarse contra GitHub. \\
\bottomrule
\end{tabular}
\end{center}

\newpage

% ============================================================
\section{Diagrama de arquitectura}

\begin{center}
\begin{tikzpicture}[
  node distance=1.1cm and 1.4cm,
  every node/.style={align=center}
]
  % Webmaster
  \node[nodoUsuario] (webmaster) {Webmaster\\\footnotesize local};

  % VS Code / navegador
  \node[nodoPrincipal, right=of webmaster, minimum width=2.7cm] (vscode)
    {VS Code\\\footnotesize edición directa};
  \node[nodoPrincipal, above=0.4cm of vscode, minimum width=2.7cm] (sveltia)
    {Sveltia CMS\\\footnotesize /admin};

  % GitHub
  \node[nodoGithub, right=2.2cm of vscode] (github)
    {GitHub\\\footnotesize repo + Actions};

  % OAuth proxy
  \node[nodoExterno, above=0.4cm of github] (oauth)
    {OAuth Proxy\\(Cloudflare)};

  % GitHub Action
  \node[nodoGithub, below=0.4cm of github, fill=itrcGold!10] (action)
    {GitHub Action\\\footnotesize build + SSH};

  % Servidor
  \node[nodoServidor, right=2.2cm of action] (server)
    {Servidor Ubuntu\\\footnotesize nginx + TLS};

  % Usuarios finales
  \node[nodoUsuario, above=1.5cm of server, text=accentGreen, draw=accentGreen, fill=accentGreen!10] (public)
    {Ciudadanía\\\footnotesize navegador};

  % === FLECHAS ===
  % Webmaster a sus herramientas
  \draw[flecha] (webmaster) -- (vscode) node[etiqueta, midway, below] {usa};
  \draw[flecha] (webmaster) to[bend left=15] (sveltia.west);

  % Sveltia usa GitHub
  \draw[flecha, dashed] (sveltia.east) to[bend left=10] node[etiqueta, above] {OAuth} (oauth.west);
  \draw[flecha, dashed] (oauth.east) to[bend left=10] (github.north);

  % VS Code y Sveltia a GitHub
  \draw[flecha] (vscode) -- (github) node[etiqueta, midway, above] {git push};
  \draw[flecha] (sveltia.east) -- ++(0.6,0) |- (github.north west) node[etiqueta, pos=0.75, above] {commit API};

  % GitHub activa Action
  \draw[flecha] (github) -- (action) node[etiqueta, midway, right] {trigger};

  % Action despliega a servidor
  \draw[flechaVerde] (action) -- (server) node[etiqueta, midway, above, text=accentGreen] {SSH + rsync};

  % Servidor sirve al público
  \draw[flechaVerde] (server) -- (public) node[etiqueta, midway, right, text=accentGreen] {HTTPS};

  % === Contenedores ===
  \begin{pgfonlayer}{background}
    \node[rectangle, rounded corners=4pt, fill=govBlue!5, draw=govBlue!30,
          inner sep=10pt, fit=(webmaster)(sveltia)(vscode),
          label={[govBlue]above:Máquina del webmaster}] {};
    \node[rectangle, rounded corners=4pt, fill=itrcGold!5, draw=itrcGold!30,
          inner sep=10pt, fit=(github)(action)(oauth),
          label={[itrcGold]above:GitHub (nube)}] {};
    \node[rectangle, rounded corners=4pt, fill=itrcNavy!5, draw=itrcNavy!30,
          inner sep=10pt, fit=(server)(public),
          label={[itrcNavy]above:Infraestructura ITRC}] {};
  \end{pgfonlayer}
\end{tikzpicture}
\end{center}

\vspace{0.5cm}
\begin{callout}[title=Cómo leer el diagrama]
Las flechas \textcolor{itrcNavy}{\textbf{azules continuas}} representan acciones del usuario o del sistema. Las \textcolor{itrcNavy}{\textbf{azules punteadas}} representan handshakes de autenticación. Las flechas \textcolor{accentGreen}{\textbf{verdes}} representan el camino final del contenido publicado hasta la ciudadanía.
\end{callout}

\newpage

% ============================================================
\section{Flujo de publicación paso a paso}

\subsection{Escenario 1: Webmaster edita con VS Code}

Aplica cuando el webmaster realiza cambios estructurales o bulk edits que no son prácticos desde el CMS.

\begin{enumerate}[leftmargin=*]
  \item El webmaster ejecuta \texttt{git pull} en su copia local para sincronizarse con los cambios más recientes.
  \item Abre el proyecto en VS Code y modifica archivos: JSON de contenido, imágenes en \texttt{public/}, o código del portal.
  \item Ejecuta \texttt{npm run dev} localmente para previsualizar cambios en \texttt{http://localhost:4321}.
  \item Satisfecho con el resultado, ejecuta \texttt{git commit} y \texttt{git push}.
  \item GitHub recibe el push y dispara automáticamente el workflow de GitHub Actions.
  \item El Action construye el sitio (\texttt{npm run build}) y sube los archivos al servidor vía SSH+rsync.
  \item nginx comienza a servir la nueva versión. Tiempo total: 2 a 4 minutos.
\end{enumerate}

\subsection{Escenario 2: Webmaster edita con Sveltia CMS}

Aplica para el flujo cotidiano: publicar noticias, actualizar documentos, cambiar banners, modificar textos.

\begin{enumerate}[leftmargin=*]
  \item El webmaster navega a \texttt{https://portal.itrc.gov.co/admin} en su navegador.
  \item Hace clic en \textit{Login with GitHub}. El flujo OAuth lo lleva a autenticarse con su cuenta.
  \item Una vez autenticado, ve el panel del CMS con las colecciones (noticias, páginas, slider, etc.).
  \item Realiza los cambios necesarios usando el editor visual.
  \item Hace clic en \textit{Publish}. Sveltia CMS crea automáticamente un commit en el repo con la autoría de su cuenta GitHub.
  \item A partir del paso 5 del escenario 1: GitHub Action dispara, construye, despliega.
\end{enumerate}

\begin{callout}[title=Trazabilidad]
Cada commit (ya sea hecho por VS Code o por Sveltia CMS) queda registrado en el historial Git con autor, fecha, y descripción. Esto permite auditar quién publicó qué y cuándo, sin necesidad de sistemas externos.
\end{callout}

\newpage

% ============================================================
\section{Configuración del servidor (Ubuntu + nginx)}

\subsection{Requisitos mínimos del servidor}

\begin{center}
\begin{tabular}{@{}ll@{}}
\toprule
\textbf{Componente} & \textbf{Versión sugerida} \\
\midrule
Sistema operativo & Ubuntu Server 22.04 LTS o superior \\
nginx & 1.18 o superior \\
Espacio en disco & Mínimo 10 GB (actual: $\sim$4 GB, con margen) \\
RAM & Mínimo 1 GB (el sitio no requiere más que esto) \\
CPU & Cualquier CPU moderno (el servidor solo sirve estáticos) \\
Puertos & 80 y 443 abiertos al público; 22 solo para administración \\
\bottomrule
\end{tabular}
\end{center}

\subsection{Usuario de despliegue}

Se crea un usuario específico para los despliegues automatizados, sin privilegios \texttt{sudo}:

\begin{verbatim}
sudo useradd -m -s /bin/bash portal-deploy
sudo mkdir -p /var/www/portal.itrc.gov.co
sudo chown portal-deploy:www-data /var/www/portal.itrc.gov.co
sudo chmod 755 /var/www/portal.itrc.gov.co
\end{verbatim}

La clave pública SSH del workflow de GitHub Actions se añade a\\\texttt{/home/portal-deploy/.ssh/authorized\_keys}.

\subsection{Configuración nginx mínima}

Archivo \texttt{/etc/nginx/sites-available/portal.itrc.gov.co}:

\begin{small}
\begin{verbatim}
server {
    listen 80;
    server_name portal.itrc.gov.co;
    return 301 https://$host$request_uri;
}

server {
    listen 443 ssl http2;
    server_name portal.itrc.gov.co;

    root /var/www/portal.itrc.gov.co;
    index index.html;

    ssl_certificate     /etc/letsencrypt/live/portal.itrc.gov.co/fullchain.pem;
    ssl_certificate_key /etc/letsencrypt/live/portal.itrc.gov.co/privkey.pem;

    # Headers de seguridad
    add_header X-Frame-Options "SAMEORIGIN" always;
    add_header X-Content-Type-Options "nosniff" always;
    add_header Referrer-Policy "strict-origin-when-cross-origin" always;
    add_header Strict-Transport-Security "max-age=31536000" always;

    # Compresión
    gzip on;
    gzip_types text/plain text/css application/javascript application/json;

    # Cache para assets estáticos
    location /_astro/ {
        expires 1y;
        add_header Cache-Control "public, immutable";
    }
    location /documentos/ {
        expires 7d;
        add_header Cache-Control "public";
    }

    # Rutas del portal
    location / {
        try_files $uri $uri.html $uri/index.html =404;
    }
}
\end{verbatim}
\end{small}

El certificado TLS se obtiene gratuitamente con Let's Encrypt vía \texttt{certbot}.

\newpage

% ============================================================
\section{Autenticación y administración de usuarios}

\subsection{Modelo de cuentas}

Cada webmaster de la institución cuenta con una \textbf{cuenta personal de GitHub}. No se comparten credenciales entre personas. Las razones son:

\begin{itemize}[leftmargin=*]
  \item \textbf{Trazabilidad}: todos los commits en el historial quedan asociados a un nombre real.
  \item \textbf{Seguridad}: si una cuenta se compromete, se revoca su acceso individualmente sin afectar al resto.
  \item \textbf{Escalabilidad}: añadir o retirar personas es cuestión de un clic, sin pasar por el equipo técnico cada vez.
  \item \textbf{Auditoría}: cualquier auditoría posterior puede reconstruir quién hizo qué con precisión.
\end{itemize}

\subsection{Flujo de autenticación (OAuth con GitHub)}

\begin{center}
\begin{tikzpicture}[
  node distance=0.9cm,
  every node/.style={align=center, font=\footnotesize}
]
  \node[nodoUsuario, minimum width=2.3cm] (u) {Webmaster};
  \node[nodoPrincipal, right=1.4cm of u, minimum width=2.6cm] (cms) {Sveltia CMS\\(/admin)};
  \node[nodoExterno, right=1.4cm of cms, minimum width=2.6cm] (proxy) {OAuth Proxy};
  \node[nodoGithub, right=1.4cm of proxy, minimum width=2.6cm] (gh) {GitHub};

  \draw[flecha] (u) -- (cms) node[etiqueta, midway, above] {1};
  \draw[flecha] (cms) -- (proxy) node[etiqueta, midway, above] {2};
  \draw[flecha] (proxy) -- (gh) node[etiqueta, midway, above] {3};
  \draw[flecha, dashed, bend right=20] (gh.south) to node[etiqueta, midway, below] {4. token} (cms.south);
\end{tikzpicture}
\end{center}

\begin{enumerate}[leftmargin=*]
  \item El webmaster entra al CMS y solicita login.
  \item El CMS redirige al OAuth Proxy (aloja credenciales de la app OAuth registrada en GitHub).
  \item El Proxy redirige a GitHub para que el usuario autorice.
  \item GitHub entrega un token al CMS via el Proxy. El CMS queda autenticado durante la sesión (típicamente 8 horas).
\end{enumerate}

\subsection{Alta de un webmaster nuevo}

El procedimiento, a cargo del coordinador técnico:

\begin{enumerate}[leftmargin=*]
  \item El webmaster crea (o ya tiene) una cuenta personal en \href{https://github.com}{github.com}. Se recomienda usar el correo institucional.
  \item El coordinador accede al repositorio del portal en GitHub, en la sección \textit{Settings → Collaborators}, y añade al nuevo webmaster con el rol \textbf{Write}.
  \item GitHub envía una invitación por correo. El webmaster la acepta.
  \item El webmaster prueba entrar a \texttt{https://portal.itrc.gov.co/admin} y hacer login. Si llega al dashboard, está listo.
\end{enumerate}

\subsection{Roles sugeridos}

\begin{center}
\begin{tabular}{@{}lll@{}}
\toprule
\textbf{Rol GitHub} & \textbf{Puede hacer} & \textbf{Para quién} \\
\midrule
\textit{Admin} & Todo + gestionar colaboradores & 1–2 coordinadores técnicos \\
\textit{Maintain} & Push + gestionar issues & Coordinadores de contenido \\
\textit{Write} & Push a main (suficiente para CMS) & Webmasters regulares \\
\textit{Read} & Solo lectura & Auditoría, visores externos \\
\bottomrule
\end{tabular}
\end{center}

\subsection{Baja de un webmaster}

Al retirar a una persona:
\begin{itemize}[leftmargin=*]
  \item Se remueve su acceso en \textit{Settings → Collaborators → Remove}.
  \item Los commits previos quedan en el historial (no se eliminan, por auditoría).
  \item No se revoca la aplicación OAuth: está asociada al portal, no a la persona.
\end{itemize}

\subsection{Buenas prácticas de seguridad}

\begin{itemize}[leftmargin=*]
  \item Exigir \textbf{autenticación en dos pasos (2FA)} en todas las cuentas que tengan acceso al repo.
  \item Revisar trimestralmente la lista de colaboradores y retirar accesos inactivos.
  \item El \textit{Client Secret} del OAuth App nunca debe aparecer en código ni en commits.
  \item Ante sospecha de cuenta comprometida: remover acceso inmediato, cambio de contraseña y rotación de 2FA por parte del usuario, y reincorporación controlada.
\end{itemize}

\newpage

% ============================================================
\section{Consideraciones operativas}

\subsection{Manejo de binarios}

Los archivos binarios del portal (PDFs, hojas de cálculo, imágenes, audios, videos) residen en el directorio \texttt{public/documentos/} del repositorio y están \textbf{versionados junto al código}. Esto implica:

\begin{itemize}[leftmargin=*]
  \item Cualquier archivo subido por Sveltia CMS se commitea automáticamente al repo.
  \item El clon inicial del repositorio es del orden de $\sim$3.5 GB. Esto es una sola vez; los \texttt{pull} posteriores solo transmiten cambios.
  \item La ventaja: una única fuente de verdad para contenido y binarios.
  \item La ruta de servido se controla por la variable \texttt{DOCS\_BASE\_URL} en el código. Para despliegue en servidor propio, el valor es \texttt{/documentos}.
\end{itemize}

\subsection{Rollback (revertir un cambio)}

Si se publica un cambio erróneo:
\begin{enumerate}[leftmargin=*]
  \item Identificar el commit problemático con \texttt{git log --oneline}.
  \item Ejecutar \texttt{git revert <hash>} — esto crea un nuevo commit que deshace los cambios.
  \item \texttt{git push}. El GitHub Action despliega la versión restaurada en 2–4 minutos.
\end{enumerate}

\begin{warnbox}[title=Nunca forzar el historial]
\textbf{No usar \texttt{git push --force}} sobre la rama principal. El historial es auditable y debe preservarse. La forma correcta de deshacer cambios publicados es siempre con \texttt{git revert}, que deja evidencia de la corrección en el log.
\end{warnbox}

\subsection{Monitoreo y logs}

\begin{itemize}[leftmargin=*]
  \item \textbf{Logs de nginx}: \texttt{/var/log/nginx/access.log} y \texttt{error.log} en el servidor. Rotación semanal estándar.
  \item \textbf{Estado de despliegues}: pestaña \textit{Actions} del repositorio en GitHub, cada workflow se registra con éxito/fallo.
  \item \textbf{Disponibilidad}: monitoreo externo opcional (uptime pings) hacia \texttt{portal.itrc.gov.co}.
\end{itemize}

\subsection{Respaldos}

El modelo provee respaldo implícito:
\begin{itemize}[leftmargin=*]
  \item El repositorio en GitHub es un backup por sí mismo (replicado en infraestructura GitHub).
  \item Cada webmaster con su clon local tiene otra copia completa.
  \item El historial Git permite restaurar cualquier versión anterior por fecha o commit.
  \item Para infraestructura crítica, el coordinador puede configurar adicionalmente un \texttt{git clone --mirror} periódico en un servidor institucional secundario.
\end{itemize}

\vspace{0.5cm}

% ============================================================
\section*{Referencias}
\addcontentsline{toc}{section}{Referencias}

\begin{itemize}[leftmargin=*]
  \item Astro: \url{https://astro.build}
  \item Sveltia CMS: \url{https://github.com/sveltia/sveltia-cms}
  \item nginx documentation: \url{https://nginx.org/en/docs/}
  \item Let's Encrypt / Certbot: \url{https://certbot.eff.org}
  \item GitHub Actions: \url{https://docs.github.com/actions}
  \item Manual del operador (portal ITRC): \texttt{docs/manual-operador/README.md} en el repo
\end{itemize}

\vfill
\begin{center}
  \color{itrcGold}\rule{0.5\linewidth}{1pt}\\[0.2cm]
  {\small\color{midGray} Documento generado como parte del proceso de migración del portal ITRC.\\ Ante dudas técnicas, consultar al equipo de TI del ITRC.}
\end{center}

\end{document}
